\section{Results}
\label{sec:results}

\subsection{Global Planetary Magnetic Topology: Moon vs. Mars}
Our comparative analysis reveals fundamental differences in the spatial structure, magnitude, and global organization of remnant crustal magnetization between the Moon and Mars (Table~\ref{tab:comparative_landing_sites}).

At the standard orbital mapping altitude of \SI{30}{\kilo\meter}, the global mean lunar crustal magnetic field magnitude is $|\mathbf{B}| = \SI{0.77}{\nano\tesla}$ (median = \SI{0.56}{\nano\tesla}), with localized peak anomalies reaching $\sim\SI{10}{\nano\tesla}$ to \SI{20}{\nano\tesla} over the South Pole-Aitken (SPA) basin and lunar impact antipodes \citep{Hood2020, Tsunakawa2015}. The overall lunar magnetic landscape is characterized by widespread, homogeneous near-null magnetic field conditions across its near-side maria and polar regions.

In contrast, Mars exhibits an unprecedented planetary magnetic dichotomy at its \SI{400}{\kilo\meter} mapping altitude (global mean $|\mathbf{B}| = \SI{15.05}{\nano\tesla}$, median = \SI{1.51}{\nano\tesla}, peak $> \SI{360}{\nano\tesla}$) \citep{Langlais2019}. The ancient Noachian crust of the southern cratered highlands ($\SI{150}{\degree}\text{E}$ to $\SI{240}{\degree}\text{E}$, $\SI{30}{\degree}\text{S}$ to $\SI{75}{\degree}\text{S}$; Terra Sirenum and Terra Cimmeria) harbors coherent, multipolar magnetic lineations extending thousands of kilometers, generating ground-level surface fields exceeding \SI{1500}{\nano\tesla}. Conversely, the northern lowlands (Vastitas Borealis, Utopia Planitia, Arcadia Planitia) and major impact basins (Hellas, Isidis, Argyre) are severely demagnetized, exhibiting orbital fields $|\mathbf{B}_{\text{400km}}| < \SI{2.0}{\nano\tesla}$ and surface fields $|\mathbf{B}_{\text{surf}}| < \SI{25}{\nano\tesla}$.

\subsection{Comparative Landing Site Characterization across 36 Planetary Targets}
Extraction of localized vector magnetic field components across 36 landing sites establishes a continuous magnetic attenuation spectrum spanning five orders of magnitude relative to Earth's geomagnetic baseline ($50\,\mu\text{T}$) (Table~\ref{tab:comparative_landing_sites}).

\subsubsection{Lunar Artemis III Candidate South Polar Regions}
All 13 candidate landing regions prioritized for the Artemis III mission ($\SI{84}{\degree}\text{S}$ to $\SI{90}{\degree}\text{S}$) reside in an ultra-low, near-null hypomagnetic field regime. Field intensities at \SI{30}{\kilo\meter} altitude range from \SI{0.40}{\nano\tesla} at Haworth crater to \SI{0.82}{\nano\tesla} at Peak Near Shackleton. All 13 candidate sites are classified as \textbf{low/null-field} environments ($< \SI{1.0}{\nano\tesla}$), representing a $>50,000$-fold to $>125,000$-fold reduction relative to Earth's GMF ($50\,\mu\text{T}$).

\subsubsection{Historical and Contemporary Lunar Landing Sites}
Historical Apollo mare sites (Apollo 11, 12, 14, 15) and recent robotic landers (Chang'e 3, 5; Chandrayaan-3) exhibit near-null fields ($0.10\text{--}\SI{1.00}{\nano\tesla}$). However, Apollo 16 in the Descartes Highlands represents a \textbf{moderate field} regime ($|\mathbf{B}_{\text{30km}}| = \SI{4.94}{\nano\tesla}$), situated adjacent to a strong highland magnetic anomaly. Chang'e 4 in the Von K\'{a}rm\'{a}n crater on the lunar far side registers \SI{0.72}{\nano\tesla} at \SI{30}{\kilo\meter} altitude within the margin of the SPA basin anomaly.

\subsubsection{Martian ISRU Candidate Habitats and Rover Sites}
Proposed human landing sites selected for shallow subsurface glacial ice access in the northern plains—Arcadia Planitia ($\SI{39.30}{\degree}\text{N}$, $\SI{189.70}{\degree}\text{E}$) and Deuteronilus Mensae ($\SI{39.10}{\degree}\text{N}$, $\SI{23.20}{\degree}\text{E}$)—reside in extreme low/null-field environments ($|\mathbf{B}_{\text{400km}}| \approx 0.70\text{--}\SI{0.71}{\nano\tesla}$, ground surface $|\mathbf{B}_{\text{surf}}| \approx 5.9\text{--}\SI{6.0}{\nano\tesla}$). Active and historical rover sites (Perseverance in Jezero Crater: surface $|\mathbf{B}| \approx \SI{23.4}{\nano\tesla}$; Zhurong in Utopia Planitia: surface $|\mathbf{B}| \approx \SI{6.5}{\nano\tesla}$; ExoMars Oxia Planum: surface $|\mathbf{B}| \approx \SI{15.5}{\nano\tesla}$) are similarly situated in severely hypomagnetic regimes ($>2,000\times$ to $>15,000\times$ lower than Earth).

\subsubsection{InSight Ground Truth and Southern Anomaly Mini-Magnetospheres}
NASA's InSight lander in Elysium Planitia ($\SI{4.50}{\degree}\text{N}$, $\SI{135.62}{\degree}\text{E}$) recorded an unexpectedly high static ground surface field of $\sim\SI{2013}{\nano\tesla}$ \citep{Johnson2020}, demonstrating that local shallow magnetized crustal basement bodies can amplify ground-level fields by an order of magnitude over orbital predictions. In the southern highlands, candidate scientific outposts in the Terra Sirenum anomaly belt ($\SI{30.00}{\degree}\text{S}$, $\SI{195.00}{\degree}\text{E}$) exhibit orbital fields of \SI{140.07}{\nano\tesla} and surface fields exceeding \SI{1190}{\nano\tesla} (regional peaks $> \SI{1500}{\nano\tesla}$), forming mini-magnetospheres capable of deflecting low-energy solar wind ions \citep{Mitchell2001}.

\begin{table*}[t]
\centering
\footnotesize
\caption{\textbf{Comparative magnetic field environment across 36 lunar and Martian landing sites.} Magnetic field magnitude at orbital mapping altitude ($|\mathbf{B}_{\text{orb}}|$: 30\,km for Moon, 400\,km for Mars) and estimated/measured ground surface intensity ($|\mathbf{B}_{\text{surf}}|$) derived from satellite magnetometry (Lunar Prospector, Kaguya, MGS, MAVEN) and InSight ground truth \citep{Hood2020, Hood2022, Langlais2019, Johnson2020}.}
\label{tab:comparative_landing_sites}
\begin{tabular}{llrrcll}
\toprule
\textbf{Landing Site / Target} & \textbf{Body / Mission} & \textbf{Lat ($^\circ$)} & \textbf{Lon ($^\circ$E)} & \textbf{$|\mathbf{B}_{\text{orb}}|$ (nT)} & \textbf{$|\mathbf{B}_{\text{surf}}|$ (nT)} & \textbf{Classification} \\
\midrule
\multicolumn{7}{l}{\textbf{\textit{Lunar Landing Sites (30\,km Altitude Baseline)}}} \\
Faustini Rim A & Moon / Artemis III & -85.30 & 77.00 & 0.45 & 0.45 & low/null-field \\
Peak Near Shackleton & Moon / Artemis III & -89.70 & 166.00 & 0.82 & 0.82 & low/null-field \\
Connecting Ridge & Moon / Artemis III & -88.00 & 137.00 & 0.60 & 0.60 & low/null-field \\
Connecting Ridge Ext. & Moon / Artemis III & -88.30 & 148.00 & 0.65 & 0.65 & low/null-field \\
de Gerlache Rim 1 & Moon / Artemis III & -88.50 & 290.00 & 0.50 & 0.50 & low/null-field \\
de Gerlache Rim 2 & Moon / Artemis III & -88.70 & 294.00 & 0.48 & 0.48 & low/null-field \\
de Gerlache-Kocher Massif & Moon / Artemis III & -86.00 & 285.00 & 0.52 & 0.52 & low/null-field \\
Haworth & Moon / Artemis III & -87.40 & 358.00 & 0.40 & 0.40 & low/null-field \\
Malapert Massif & Moon / Artemis III & -86.00 & 0.00 & 0.55 & 0.55 & low/null-field \\
Leibnitz Beta Plateau & Moon / Artemis III & -85.00 & 31.00 & 0.70 & 0.70 & low/null-field \\
Nobile Rim 1 & Moon / Artemis III & -85.40 & 35.00 & 0.62 & 0.62 & low/null-field \\
Nobile Rim 2 & Moon / Artemis III & -84.80 & 47.00 & 0.58 & 0.58 & low/null-field \\
Amundsen Rim & Moon / Artemis III & -84.50 & 85.00 & 0.65 & 0.65 & low/null-field \\
Apollo 11 & Moon / Apollo & 0.67 & 23.47 & 0.93 & 0.93 & low/null-field \\
Apollo 12 & Moon / Apollo & -3.01 & 336.58 & 1.00 & 1.00 & low/null-field \\
Apollo 14 & Moon / Apollo & -3.65 & 342.53 & 0.44 & 0.44 & low/null-field \\
Apollo 15 & Moon / Apollo & 26.13 & 3.63 & 0.23 & 0.23 & low/null-field \\
Apollo 16 (Highlands) & Moon / Apollo & -8.97 & 15.50 & 4.94 & 4.94 & moderate-field \\
Apollo 17 & Moon / Apollo & 20.19 & 30.77 & 1.08 & 1.08 & moderate-field \\
Chang'e 3 & Moon / Chang'e & 44.12 & 340.49 & 0.19 & 0.19 & low/null-field \\
Chang'e 4 (SPA Basin) & Moon / Chang'e & -45.46 & 177.60 & 0.72 & 0.72 & low/null-field \\
Chang'e 5 & Moon / Chang'e & 43.06 & 308.08 & 0.10 & 0.10 & low/null-field \\
Chandrayaan-3 & Moon / Chandrayaan & -69.37 & 32.35 & 0.78 & 0.78 & low/null-field \\
\midrule
\multicolumn{7}{l}{\textbf{\textit{Martian Landing Sites (400\,km Altitude Baseline \& Surface Modeling)}}} \\
Jezero Crater (Perseverance) & Mars / Mars 2020 & 18.45 & 77.45 & 2.75 & 23.4 & low/null-field \\
Gale Crater (Curiosity) & Mars / MSL & -4.59 & 137.44 & 3.18 & 27.0 & low/null-field \\
Elysium Planitia (InSight) & Mars / InSight & 4.50 & 135.62 & 1.99 & 2013.0 & high-field outpost \\
Meridiani Planum (Opportunity) & Mars / MER-B & -1.95 & 354.47 & 2.55 & 21.6 & low/null-field \\
Gusev Crater (Spirit) & Mars / MER-A & -14.57 & 175.47 & 9.08 & 77.2 & low/null-field \\
Vastitas Borealis (Phoenix) & Mars / Scout & 68.22 & 234.25 & 0.48 & 4.1 & low/null-field \\
Chryse Planitia (Viking 1) & Mars / Viking & 22.48 & 312.05 & 0.59 & 5.0 & low/null-field \\
Utopia Planitia (Viking 2) & Mars / Viking & 47.97 & 134.28 & 0.39 & 3.3 & low/null-field \\
Utopia Planitia (Zhurong) & Mars / Tianwen-1 & 25.07 & 109.92 & 0.76 & 6.5 & low/null-field \\
Oxia Planum (ExoMars) & Mars / Rosalind Franklin & 18.27 & 335.37 & 1.82 & 15.5 & low/null-field \\
Arcadia Planitia (Human Candidate) & Mars / Human ISRU Base & 39.30 & 189.70 & 0.71 & 6.0 & low/null-field \\
Deuteronilus Mensae (Candidate) & Mars / Human ISRU Base & 39.10 & 23.20 & 0.70 & 5.9 & low/null-field \\
Terra Sirenum Anomaly Belt & Mars / Science Outpost & -30.00 & 195.00 & 140.07 & 1190.6 & high-field outpost \\
\bottomrule
\end{tabular}
\end{table*}


\subsection{Systematic Review of Hypomagnetic Literature}
Our systematic review (Table~\ref{tab:unified_biology_literature}) documents consistent physiological perturbations across phylogenetic groups exposed to HMF conditions ($<\SI{5}{\micro\tesla}$):

\begin{enumerate}
    \item \textbf{Plant Morphogenesis and Physiology:} In \textit{Arabidopsis thaliana}, near-null magnetic fields ($<\SI{1}{\micro\tesla}$) delay floral transition via down-regulation of cryptochrome (CRY1/CRY2) and phytochrome photoperiodic signaling pathways \citep{Agliassa2018}. HMF conditions disrupt root gravitropic curvature and meristem cell elongation by altering PIN2 polar auxin transport \citep{Narayana2018}, downregulate iron uptake transcription factors (FIT, IRT1) \citep{Narayana2021}, and elevate reactive oxygen species (ROS) \citep{Maffei2014, Belyavskaya2004}.
    \item \textbf{Microbial Ecology and Magnetotaxis:} Bacterial cultures in HMFs exhibit altered cell division rates and membrane transport kinetics \citep{Creanga2009}, phylum-level shifts in gut microbiome community composition \citep{Zhang2017}, and complete loss of magnetotaxis in magnetotactic species \citep{Frankel1981}.
    \item \textbf{Animal and Human Cellular Systems:} In invertebrate models (\textit{Drosophila}, \textit{C. elegans}), HMF exposure extends circadian period lengths via radical pair spin state recombination modulation \citep{Yoshii2009} and alters sensory navigation \citep{VidalGadea2015}. In mammalian and human cell models, HMF exposure alters neuroblastoma cell cycle transcriptomes \citep{Mo2012}, modifies microcirculatory capillary blood flow \citep{Gurfinkel2014}, and acts synergistically with mechanical unloading to accelerate trabecular bone demineralization and osteoclast resorption \citep{Xu2012, Mo2014}.
\end{enumerate}

\input{../tables/table2_unified_biology_literature.tex}

\subsection{Multi-Planetary Biological Risk Matrix}
Synthesizing physical field determinations with empirical biological sensitivity thresholds yields the multi-planetary risk assessment matrix presented in Table~\ref{tab:unified_risk_matrix}.

\begin{table*}[t]
\centering
\scriptsize
\caption{\textbf{Comparative biological risk assessment matrix across lunar and Martian surface magnetic field regimes.} Vulnerability of critical biological systems relative to terrestrial geomagnetic field baseline ($50\,\mu\text{T}$).}
\label{tab:unified_risk_matrix}
\begin{tabular}{p{3.2cm}cccccc}
\toprule
\textbf{Biological System} & \textbf{Earth GMF} & \textbf{Mars Sirenum} & \textbf{Mars InSight} & \textbf{Mars Northern Plains} & \textbf{Lunar Highlands} & \textbf{Lunar South Pole} \\
 & ($\sim\SI{50}{\micro\tesla}$) & ($>\SI{1000}{\nano\tesla}$) & ($\sim\SI{2000}{\nano\tesla}$) & ($<\SI{25}{\nano\tesla}$) & ($\SI{1}{--}\SI{5}{\nano\tesla}$) & ($<\SI{1.0}{\nano\tesla}$) \\
\midrule
Plant Vegetative Biomass & Baseline & Low Risk & Low Risk & High Risk & Moderate Risk & High Risk \\
Plant Flowering \& Seed Set & Baseline & Moderate Risk & Moderate Risk & High Risk & Moderate-High Risk & High Risk \\
Closed-Loop Microbiome & Baseline & Low Risk & Low Risk & Moderate Risk & Low-Moderate Risk & Moderate Risk \\
Circadian Clocks & Baseline & Low-Moderate Risk & Low-Moderate Risk & High Risk & Moderate-High Risk & High Risk \\
Human Musculoskeletal & Baseline & Moderate Risk & Moderate Risk & High Risk & High Risk & High Risk \\
Cellular DNA / ROS Balance & Baseline & Low-Moderate Risk & Low-Moderate Risk & High Risk & Moderate-High Risk & High Risk \\
\bottomrule
\end{tabular}
\end{table*}

